\documentclass[11pt]{article}
\usepackage[margin=1in]{geometry}
\usepackage[T1]{fontenc}
\usepackage{lmodern,amsmath,amssymb,booktabs,graphicx,enumitem}
\usepackage[round,authoryear]{natbib}
\usepackage[colorlinks=true,linkcolor=blue,urlcolor=blue,citecolor=blue]{hyperref}
\setlength{\parindent}{0pt}
\setlength{\parskip}{5pt}
\setlength{\emergencystretch}{2em}
\title{RoCE-K: fallback, calibration and nuisance-error experiments}
\author{}
\date{2 October 2026}
\begin{document}
\maketitle
\textbf{Status.} This study tests repairs to the residual-compatibility
prototype. It includes 18 Gaussian settings with 2,000 repetitions each,
18 binary-score settings with 1,000 each, and 12 fresh saturated-nuisance
settings with 300 each: 57,600 setting-specific repetitions. A paired
refinement of the nuisance budgets reuses the same observations and is not
counted again. All reported sites contain 1,000 patients. Separately, 90
archived calibrated RoCE fits and 40,000 target-only reference repetitions
are audited. No production estimator or historical result is replaced.

\section{The repairs and their statistical roles}

Retain $\tau=d+r_1-r_0$. Each source arm provides a Fourier confidence set
intersected with a dispersion interval. The target supplies a joint region
for $(d,r_1,r_0)$. An allocated union bound protects projection onto TATE;
the two arms and the target summaries need not be independent.

\paragraph{A pointwise length-controlled fallback.}
Let $J_1,J_0$ be the protected source dispersion intervals and
$C_a\subseteq J_a$ their intersections with Fourier constraints. If the
target region projects onto $d\in[\widehat d-u_d,\widehat d+u_d]$, define
\[
 J_S=[\widehat d-u_d+\inf J_1-\sup J_0,\;
       \widehat d+u_d+\sup J_1-\inf J_0].
\]
When the full target/source intersection is nonempty, report its projection.
When it is empty, construct a source interval using the hulls of $C_1,C_0$
if both exist, and otherwise using $J_1,J_0$. Return the shorter of this
source interval and the prespecified target reference.

On the simultaneous coverage event the full intersection cannot be empty,
so changing the fallback does not remove that event. Moreover, for every
realized dataset the new reported interval has length at most $|J_S|$.
This removes the separate target-width fallback contribution from an
expected-length bound. A useful bound on $E|J_S|$ still requires its own
conditions, especially with fitted nuisances and estimated variances.

The fixed allocation uses target failure probability $.025$. The decreasing
alternative uses $.025/\sqrt K$. After reserving any variance and bias
certificate errors, the remaining allowance is split between arms, and
within each arm between Fourier and dispersion. These rules are specified
before the new evaluations. We compare the interval midpoint with a second
point rule: form a source-compatible residual center, add $\widehat d$, and
project that point into the final interval. Neither rule is asserted to be
MSE-optimal.

\paragraph{Sharper target regions.}
For bounded patient scores, use simultaneous bands in four directions:
the three coordinates and the TATE contrast $(1,1,-1)$. The latter uses the
full score covariance. The target region is the intersection of those bands;
projection with residual interval sets is exact interval arithmetic.

Known-variance Bernstein is an oracle comparison. The empirical version
uses the unbiased sample variance and the deterministic score range, via
Theorem 4 of \citet{maurer2009}, with explicit two-sided and four-direction
error allocation. Source variance certificates use that paper's Theorem 10
on an independent pilot, unioned over sites and arms. They are confidence
bounds, not empirical variances declared to be true. Non-Gaussian source
characteristic functions retain the bounded-score approximation allowance
from the preceding experiment.

\section{Gaussian results: fallback repair works}

The table uses $K=2048$, no shared target variance floor, and 2,000 new-seed
repetitions. Length is relative to the ordinary target-only normal interval.
The Gaussian experiment is an abstract joint residual model; it is not
asserted to be an exact binary causal-data embedding.

\begin{center}\small
\resizebox{\linewidth}{!}{\input{gaussian_table}}
\end{center}

The shorter fallback substantially improves midpoint RMSE without discarding
any failed intersection. At the all-valid setting its RMSE falls from about
$.0074$ to $.0010$. At strong shifts it falls from about $.0076$ to $.0021$.
The source-compatible point is also recorded, and can improve on the midpoint
in some cases. It is not uniformly preferable.

Decreasing the target error allocation helps when the shared target component
has zero variance, but can increase length when that component is noisy.
At the weak-shift setting with shared target SD equal to half a source SE,
the fixed shorter interval has length about $.526$ of target-only, versus
$.695$ for the decreasing allocation. The fixed shorter rule is therefore
the preferred candidate for subsequent work; decreasing allocation is retained
as an explicit alternative, not selected after seeing the interval.

\begin{figure}[p]
\centering\includegraphics[width=\linewidth]{fallback_comparison.pdf}
\caption{Paired Gaussian comparison at weak shifts. The source-based fallback
controls length on every repetition. All empty-intersection cases remain.}
\end{figure}

\section{Binary scores: calibration improves, variance uncertainty remains}

The binary model retains randomized treatment, binary covariates and oracle
common outcome scores from the preceding study. Target truth is $.10$.
Sources now use 250 variance-pilot observations and 750 evaluation observations;
the target uses 1,000. This respects the total site budget.

At $K=2048$, weak shifts and effect heterogeneity $.3$, the length ratios
relative to ordinary target-only are approximately:
\begin{center}
\begin{tabular}{lr}
\toprule
Calibration & Length/target\\\midrule
Known variance, target moment ellipsoid & 2.056\\
Known variance, target Bernstein region & 0.924\\
Known source variance, empirical target Bernstein & 1.089\\
Estimated source variance, empirical target Bernstein & 1.737\\
Gaussian plug-in diagnostic & 0.721\\\bottomrule
\end{tabular}
\end{center}
Thus the loose target ellipsoid was an avoidable cost. However, the estimated
source-variance certificate still loses substantial precision. The short
Gaussian plug-in result is not promoted to a uniformly valid procedure on
the basis of this experiment. These panels continue to use oracle nuisance
functions, so they do not by themselves establish fully feasible inference.

\begin{figure}[p]
\centering\includegraphics[width=\linewidth]{binary_calibration.pdf}
\caption{Binary-score calibration, 1,000 repetitions per setting. Source
variance estimation uses a separate pilot within the 1,000-patient budget.}
\end{figure}

\section{Fresh nuisance fitting and feasible bias budgets}

This diagnostic uses two binary covariates and four strata. Target outcome
means are $.55+cX_1/2+.08X_2$ and $.45-cX_1/2+.08X_2$, so TATE is exactly
$.10$. Source covariate distributions and treatment propensities differ.
Some treated source arms receive a local mean shift. Each source uses
500 training, 250 variance-pilot and 250 evaluation patients. The target
uses 500 training and 500 evaluation patients.

Every repetition refits half-count-regularized saturated outcome and
propensity probabilities. Source merged weights calibrate regularized
source joint-cell frequencies to target cell frequencies. This is a
closed-form low-dimensional calibration diagnostic, not the production
sparse fold-summed RoCE fitting program.

Let $\widehat p_t(x)$ and $\widehat p_{j,a}(x)$ denote fitted target and source
joint-cell probabilities, and $\widehat q_j(x)=\widehat p_t(x)/\widehat
p_{j,a}(x)$. For a structurally valid source, its conditional nuisance drift
equals
\[
 \sum_x\{p_{j,a}(x)\widehat q_j(x)-p_t(x)\}
                     \{m_t^a(x)-\widehat m^a(x)\}.
\]
Simultaneous binomial confidence bounds give a computable allowance
\[
 \sum_x\{\widehat q_j(x)\Delta p_{j,a}(x)+\Delta p_t(x)\}
                  \Delta m_a(x).
\]
The exact conditional drift is also computed from the known simulation law,
solely as an oracle diagnostic. It is never mislabeled as a feasible bound.

The first allocation distributes the training confidence error over all
coordinates. A paired refinement assigns half to the shared target nuisance
events once, independently of $K$, and half to source events. At $K=128$,
typical maximum oracle source budgets are about $.014$, while the original
feasible budgets are about $.34$ and the shared-target refinement about
$.28$. The latter is an improvement, but still far too large for useful
borrowing in this diagnostic.

The ordinary target reference uses all 1,000 target patients in a two-fold
AIPW estimator. Its score-normal interval is an asymptotic comparator.
An independent 20,000-repetition audit per heterogeneity setting finds
coverage $.94915$ and $.94535$, with Monte Carlo standard errors about $.0016$.
We also construct a fully protected target-only reference: each fold has
training-error allowance $.0025$ and held-out mean-error allowance $.0225$;
averaging its two protected intervals is justified by a union bound.

With feasible shared-target budgets, the multisource interval is about
$8.7$--$8.8$ times the ordinary normal reference and about $1.10$--$1.11$
times this similarly protected target-only interval. Therefore this feasible
finite-sample implementation is not yet an efficient borrowing method.
Its high coverage reflects conservatism, not a completed high-dimensional
solution. The paired refinement reproduces the first five original method
outputs exactly and is not counted as additional independent evidence.

\begin{figure}[p]
\centering\includegraphics[width=\linewidth]{fitted_cost.pdf}
\caption{Fresh saturated nuisance fits, 300 repetitions per setting. Oracle
budgets locate the cost of protection but are unavailable in actual analysis.}
\end{figure}

\section{Archived RoCE fits, checks and the next priority}

Ninety previously completed actual RoCE calibration fits are read without
modification, with all completion hashes verified. The common-outcome scores
are transformed into the target/residual coordinates. Empirical covariance
exports explicitly convert the old $n^{-2}$ centered-moment convention to
the unbiased sample-mean covariance convention. Conditional population moments
remain labeled as oracle inputs. The largest archived valid-source absolute
drift is about $.0295$, and the largest target-coordinate drift about $.00875$.
These heterogeneous single-fit observations do not constitute a new Monte
Carlo coverage experiment, but confirm that zero drift is not automatic.

Implementation checks cover analytic projection against independent linear
and nonlinear solvers, the pointwise fallback length cap, fitted-score
identities, simultaneous variance bounds and the total patient budget.
The previous binary generator reproduces exactly under its old settings.
All interval, coverage, length and point-error summaries are recomputed from
saved repetitions. Old experiments, preliminary fits and paired revisions
are preserved on scratch.

The immediate decision is to keep the fixed-budget shorter fallback and
the sharper target calibration as research components. The remaining obstacle
is the cost of simultaneous variance and nuisance-error protection. Subsequent
work should use the calibration equations to control the combined remainder
more tightly, and establish joint $(n,K)$ approximations where appropriate.
Neither ignoring the nuisance budget nor assuming estimated covariance is
known is a justified shortcut. Full sparse/high-dimensional fitting should
follow a useful uncertainty analysis, rather than precede it.

\bibliographystyle{plainnat}
\bibliography{references}
\end{document}
